% Referências em formato ABNT

% === FEDERATED LEARNING - FUNDAMENTOS ===

\bibitem{mcmahan2017fedavg}McMAHAN, B. et al.
Communication-Efficient Learning of Deep Networks from Decentralized Data.
{\bf Artificial Intelligence and Statistics (AISTATS)},
p.1273-1282, 2017.

\bibitem{yang2019fl}YANG, Q. et al.
Federated Learning.
{\bf Synthesis Lectures on Artificial Intelligence and Machine Learning},
v.13, n.3, p.1-207, 2019.

\bibitem{flower2020}BEUTEL, D.J. et al.
Flower: A Friendly Federated Learning Research Framework.
{\bf arXiv preprint arXiv:2007.14390}, 2020.

\bibitem{bonawitz2019scale}BONAWITZ, K. et al.
Towards Federated Learning at Scale: System Design.
{\bf Proceedings of Machine Learning and Systems (MLSys)},
v.1, p.374-388, 2019.

\bibitem{konecny2016communication}KONEČNÝ, J. et al.
Federated Learning: Strategies for Improving Communication Efficiency.
{\bf arXiv preprint arXiv:1610.05492}, 2016.

% === SURVEYS DE FEDERATED LEARNING ===

\bibitem{liu2024survey}LIU, B. et al.
Recent advances on federated learning: A systematic survey.
{\bf Neurocomputing},
v.597, p.128019, 2024.

\bibitem{chai2024survey}CHAI, D. et al.
A Survey for Federated Learning Evaluations: Goals and Measures.
{\bf arXiv preprint arXiv:2308.11841}, 2024.

\bibitem{shi2023survey}SHI, L.
A survey on federated learning: evolution, applications and challenges.
{\bf Proceedings of the 5th International Conference on Computing and Data Science},
DOI: 10.54254/2755-2721/22/20231177, 2023.

\bibitem{abdulrahman2021survey}ABDULRAHMAN, S. et al.
A survey on federated learning: The journey from centralized to distributed on-site learning and beyond.
{\bf IEEE Internet of Things Journal},
v.8, n.7, p.5476-5497, 2021.

\bibitem{nguyen2021iot}NGUYEN, D.C. et al.
Federated learning for internet of things: A comprehensive survey.
{\bf IEEE Communications Surveys \& Tutorials},
v.23, n.2, p.1622-1658, 2021.

% === FEDERATED LEARNING EM FINANÇAS ===

\bibitem{kennedy2025financial}KENNEDY, C.H.; HILAL, A.; MOMENI, M.
The Role of Federated Learning in Improving Financial Security: A Survey.
{\bf Proceedings of the IEEE GCAIoT 2025 Conference}, 2025.

\bibitem{farrukh2025blockchain}FARRUKH, H. et al.
Blockchain-Based Fraud Detection: A Comparative Systematic Literature Review of Federated Learning and Machine Learning Approaches.
{\bf Electronics},
v.14, p.4952, 2025.

\bibitem{claus2025collaborative}CLAUS, B.; JOHN, A.; JOHN, A.
Federated Learning for Collaborative Fraud Detection in Financial Networks: Addressing Data Privacy Concerns.
{\bf arXiv preprint}, 2025.

\bibitem{lynch2024fraud}LYNCH, S.; ABDELMONIEM, A.M.; GILL, S.S.
Centralised and Decentralised Fraud Detection Approaches in Federated Learning.
{\bf Applications of AI for Interdisciplinary Research},
Taylor \& Francis Group, 2024.

% === MODELOS BASEADOS EM ÁRVORE FEDERADOS ===

\bibitem{cheng2021secureboost}CHENG, K. et al.
SecureBoost: A lossless federated learning framework.
{\bf IEEE Intelligent Systems},
v.36, n.6, p.87-98, 2021.

\bibitem{li2023fedtree}LI, Q. et al.
FedTree: A fast, easier, and more secure federated learning system for trees.
{\bf Proceedings of the VLDB Endowment}, 2023.

\bibitem{le2021fedxgboost}LE, N.K. et al.
FedXGBoost: Privacy-Preserving XGBoost for Federated Learning.
{\bf arXiv preprint arXiv:2106.10662}, 2021.

\bibitem{li2020practical}LI, Q.; WEN, Z.; HE, B.
Practical Federated Gradient Boosting Decision Trees.
{\bf Proceedings of the AAAI Conference on Artificial Intelligence},
v.34, n.04, p.4642-4649, 2020.

\bibitem{liu2020fedforest}LIU, Y. et al.
Federated Forest.
{\bf IEEE Transactions on Big Data},
v.6, n.6, p.1223-1228, 2020.

% === MODELOS BASEADOS EM ÁRVORE (CENTRALIZADOS) ===

\bibitem{chen2016xgboost}CHEN, T.; GUESTRIN, C.
XGBoost: A Scalable Tree Boosting System.
{\bf Proceedings of the 22nd ACM SIGKDD International Conference on Knowledge Discovery and Data Mining},
p.785-794, 2016.

\bibitem{ke2017lightgbm}KE, G. et al.
LightGBM: A Highly Efficient Gradient Boosting Decision Tree.
{\bf Advances in Neural Information Processing Systems (NeurIPS)},
v.30, p.3146-3154, 2017.

\bibitem{prokhorenkova2018catboost}PROKHORENKOVA, L. et al.
CatBoost: unbiased boosting with categorical features.
{\bf Advances in Neural Information Processing Systems (NeurIPS)},
v.31, p.6638-6648, 2018.

% === HETEROGENEIDADE E NON-IID ===

\bibitem{zhu2021noniid}ZHU, H. et al.
Federated learning on non-IID data: A survey.
{\bf Neurocomputing},
v.465, p.371-390, 2021.

\bibitem{li2020fedprox}LI, T. et al.
Federated Optimization in Heterogeneous Networks.
{\bf Proceedings of Machine Learning and Systems (MLSys)},
v.2, p.429-450, 2020.

\bibitem{karimireddy2020scaffold}KARIMIREDDY, S.P. et al.
SCAFFOLD: Stochastic controlled averaging for federated learning.
{\bf International Conference on Machine Learning (ICML)},
p.5132-5143, 2020.

\bibitem{reisizadeh2020fedpaq}REISIZADEH, A. et al.
FedPAQ: A Communication-Efficient Federated Learning Method with Periodic Averaging and Quantization.
{\bf International Conference on Artificial Intelligence and Statistics (AISTATS)},
p.2021-2031, 2020.

% === FAIRNESS E VIÉS EM ML ===

\bibitem{hardt2016equality}HARDT, M.; PRICE, E.; SREBRO, N.
Equality of Opportunity in Supervised Learning.
{\bf Advances in Neural Information Processing Systems (NeurIPS)},
v.29, p.3315-3323, 2016.

\bibitem{huang2022fairness}HUANG, W. et al.
Fairness and Accuracy in Horizontal Federated Learning.
{\bf Information Sciences},
v.589, p.170-185, 2022.

\bibitem{zhang2020fairfl}ZHANG, D.Y.; KOU, Z.; WANG, D.
FairFL: A Fair Federated Learning Approach to Reducing Demographic Bias in Privacy-Sensitive Classification Models.
{\bf IEEE International Conference on Big Data (Big Data)},
p.1051-1060, 2020.

\bibitem{barocas2019fairness}BAROCAS, S.; HARDT, M.; NARAYANAN, A.
{\bf Fairness and Machine Learning}.
fairmlbook.org, 2019.

\bibitem{corbett2017fairness}CORBETT-DAVIES, S. et al.
Algorithmic Decision Making and the Cost of Fairness.
{\bf Proceedings of the 23rd ACM SIGKDD International Conference on Knowledge Discovery and Data Mining},
p.797-806, 2017.

% === PERSONALIZAÇÃO EM FL ===

\bibitem{fallah2020personalized}FALLAH, A.; MOKHTARI, A.; OZDAGLAR, A.
Personalized Federated Learning with Theoretical Guarantees: A Model-Agnostic Meta-Learning Approach.
{\bf Advances in Neural Information Processing Systems (NeurIPS)},
v.33, p.3557-3568, 2020.

% === PRIVACIDADE E SEGURANÇA ===

\bibitem{bonawitz2017secagg}BONAWITZ, K. et al.
Practical Secure Aggregation for Privacy-Preserving Machine Learning.
{\bf Proceedings of the 2017 ACM SIGSAC Conference on Computer and Communications Security},
p.1175-1191, 2017.

\bibitem{geyer2017differential}GEYER, R.C.; KLEIN, T.; NABI, M.
Differentially private federated learning: A client level perspective.
{\bf arXiv preprint arXiv:1712.07557}, 2017.

\bibitem{wei2020privacy}WEI, K. et al.
Federated Learning with Differential Privacy: Algorithms and Performance Analysis.
{\bf IEEE Transactions on Information Forensics and Security},
v.15, p.3454-3469, 2020.

\bibitem{zhu2019deepleakage}ZHU, L.; LIU, Z.; HAN, S.
Deep leakage from gradients.
{\bf Advances in Neural Information Processing Systems (NeurIPS)},
v.32, 2019.

\bibitem{miao2022byzantine}MIAO, Y. et al.
Privacy-Preserving Byzantine-Robust Federated Learning via Blockchain Systems.
{\bf IEEE Transactions on Information Forensics and Security},
v.17, p.2848-2861, 2022.

\bibitem{dwork2006differential}DWORK, C. et al.
Calibrating noise to sensitivity in private data analysis.
{\bf Theory of Cryptography Conference},
p.265-284, 2006.

% === APLICAÇÕES EM IoT E MOBILE ===

\bibitem{zhao2020iot}ZHAO, Y. et al.
Privacy-Preserving Blockchain-Based Federated Learning for IoT Devices.
{\bf IEEE Internet of Things Journal},
v.8, p.1817-1829, 2020.

\bibitem{hard2018keyboard}HARD, A. et al.
Federated learning for mobile keyboard prediction.
{\bf arXiv preprint arXiv:1811.03604}, 2018.

% === FRAMEWORKS E PLATAFORMAS ===

\bibitem{liu2021fate}LIU, Y. et al.
FATE: An Industrial Grade Platform for Collaborative Learning with Data Protection.
{\bf Journal of Machine Learning Research (JMLR)},
v.22, n.226, p.1-6, 2021.

\bibitem{flower2024whatisfl}FLOWER.
What is Federated Learning?
Disponível em: \url{https://flower.ai/docs/framework/tutorial-series-what-is-federated-learning.html}.
Acesso em: 14 mar. 2026.

% === DATASET BAF ===

\bibitem{jesus2022baf}JESUS, S. et al.
Turning the Tables: Biased, Imbalanced, Dynamic Tabular Datasets for ML Evaluation.
{\bf Advances in Neural Information Processing Systems (NeurIPS)},
Track on Datasets and Benchmarks, 2022.

\bibitem{xu2019ctgan}XU, L. et al.
Modeling Tabular Data using Conditional GAN.
{\bf Advances in Neural Information Processing Systems (NeurIPS)}, 2019.

% === DETECÇÃO DE FRAUDES ===

\bibitem{pozzolo2015calibrating}DAL POZZOLO, A. et al.
Calibrating Probability with Undersampling for Unbalanced Classification.
{\bf IEEE Symposium Series on Computational Intelligence (SSCI)},
Cape Town, South Africa, 2015.

\bibitem{lopez2016paysim}LOPEZ-ROJAS, E.A.; ELMIR, A.; AXELSSON, S.
PaySim: A Financial Mobile Money Simulator for Fraud Detection.
{\bf 28th European Modeling and Simulation Symposium (EMSS)},
Larnaca, Cyprus, 2016.

% === REGULAMENTAÇÕES DE PRIVACIDADE ===

\bibitem{gdpr2016}EUROPEAN PARLIAMENT AND COUNCIL.
Regulation (EU) 2016/679 of the European Parliament and of the Council of 27 April 2016 on the protection of natural persons with regard to the processing of personal data.
{\bf Official Journal of the European Union}, v.59, 2016.

\bibitem{voigt2017gdpr}VOIGT, P.; VON DEM BUSSCHE, A.
{\bf The EU General Data Protection Regulation (GDPR): A Practical Guide}.
Springer, 2017.
